\section*{Chapitre \ref{ch:m3:getting-access-commodities-and-power} --- Accès aux marchés~: matières premières et électricité}

\begin{solution}{exo:m3:getting-access-commodities-and-power:1}
$50 \times 8\,760 = 438\,000$~MWh~; à \qty{85}{EUR/MWh}, soit 37.23~millions d'euros.
\end{solution}

\begin{solution}{exo:m3:getting-access-commodities-and-power:2}
Une limite de position est un cap strict~; un \omterm{def:m3:getting-access-commodities-and-power:accountability}{seuil de déclaration} (accountability level) est un seuil au-delà duquel la bourse peut demander
des explications et ordonner au trader de ne pas augmenter la position.
\end{solution}

\begin{solution}{exo:m3:getting-access-commodities-and-power:3}
La catégorie 1 (ring dealing) peut négocier dans le Ring~; la catégorie 3 (compensation des transactions associées) compense ses propres transactions
et ne traite qu'avec les catégories 1, 2 et 4.
\end{solution}

\begin{solution}{exo:m3:getting-access-commodities-and-power:4}
Position nette $120 - 30 - 80 = +10$~MWh~: elle est longue de 10~MWh. Elle doit vendre 10~MWh de plus (en infrajournalier) ou réviser sa prévision
avant la date limite~; à défaut, elle sera réglée comme un déséquilibre long.
\end{solution}

\begin{solution}{exo:m3:getting-access-commodities-and-power:5}
Premier cas~: $3 + \max(-1, 0) - 2 = $ 1~million d'euros. Deuxième cas~: $3 + 0.5 - 2 = $ 1.5~million d'euros.
\end{solution}

\begin{solution}{exo:m3:getting-access-commodities-and-power:6}
Le \omterm{def:m3:getting-access-commodities-and-power:remit}{REMIT} couvre les produits énergétiques de gros, y compris les dérivés sur l'électricité et le gaz négociés sur des marchés organisés~; la société est un participant de marché tenu de s'enregistrer et de déclarer,
sauf si elle ne négocie que des produits
situés hors de son champ d'application.
\end{solution}

\begin{solution}{exo:m3:getting-access-commodities-and-power:7}
Bourse~: 4.47~millions d'euros de collatéral, 0.283~million d'euros par an~; OTC~: 6.21~millions d'euros de crédit, 0.118
million d'euros~; physique~: 0.5~million d'euros, 0.175~million d'euros. La voie bourse coûte autant que la voie OTC lorsque
$60\,000 + m \times 37.23\text{ millions d'euros} \times 5\% = 118\,075$~: un taux de marge de 3.1~\%.
\end{solution}

\begin{solution}{exo:m3:getting-access-commodities-and-power:8}
Les contrats à terme couvrent le prix, et non le crédit de l'acheteur~: si l'acheteur fait défaut après avoir pris livraison, la société perd l'énergie non payée
et doit remplacer les livraisons prévues aux prix du marché, tandis que sa couverture par contrats à terme évolue
en sens inverse. Le crédit et le prix sont des risques distincts.
\end{solution}

\begin{solution}{pb:m3:getting-access-commodities-and-power:1}
\textbf{1.} 438\,000~MWh, 37.23~millions d'euros. \textbf{2.} 4.47~millions d'euros. \textbf{3.} La livraison physique
via le réseau doit être programmée, et ses déséquilibres réglés par un \omterm{def:m3:power-markets-trading:brp}{responsable d'équilibre}. \textbf{4.} 6.21
million d'euros~: le paiement suit la livraison, et deux mois couvrent un cycle de facturation, plus le temps d'arrêter les livraisons et
de remplacer l'acheteur. \textbf{5.} L'acheteur, dont la banque émet la \omterm{def:m3:getting-access-commodities-and-power:lc}{lettre de crédit} moyennant des frais (environ 93\,000~euros par an à 1.5~\%).
\textbf{6.} La société dépose 4.97~millions d'euros (marge et garantie)~; la banque de son acheteur ajoute 6.21~millions d'euros de
crédit. \textbf{7.} 248\,380~euros à 5~\%. \textbf{8.} Le montant double à peu près, car le taux de marge s'applique à un
notionnel plus important, et la marge de variation suit le prix dans l'intervalle. \textbf{9.} Il s'agit d'une promesse
de paiement de la banque envers la société, obtenue par l'acheteur, et non de liquidités déposées par la société. \textbf{10.} Le contrat de \omterm{def:m3:getting-access-commodities-and-power:group}{groupe d'équilibrage}, un
système de comptage et de programmation, les procédures de nomination auprès de l'opérateur, et l'enregistrement auprès du régulateur.
\textbf{11.} L'enregistrement auprès d'un régulateur national, qui attribue un code ACER à faire figurer sur ses déclarations de transactions.
\textbf{12.} Non, si son activité de négoce sur matières premières est accessoire à une activité principale qui n'est ni un service d'investissement
ni une activité bancaire et remplit les critères du règlement délégué~; une société de trading pure doit le vérifier. \textbf{13.} Elle doit
désigner un représentant dans l'Union et s'enregistrer dans un État membre où elle est active. \textbf{14.} Ses
contrats de gros sur l'énergie, y compris les transactions sur bourse (généralement déclarées pour elle par la bourse) et les
opérations bilatérales physiques. \textbf{15.} Les limites de position ou \omterm{def:m3:getting-access-commodities-and-power:accountability}{seuils de déclaration} de la bourse. \textbf{16.} La
voie boursière~: le collatéral est en espèces, mais les opérations sont les plus simples~; l'équilibrage physique peut suivre une fois les systèmes
et le personnel sont en place. \textbf{17.} Liquidité~: appels de marge si les prix bougent, avant tout revenu de la
vente physique. \textbf{18.} Capital, actionnariat et direction, contrôles des risques, le plan de trading et ses limites,
et le collatéral qu'elle peut déposer. \textbf{19.} \emph{Résultat chiffré~:} 4.97~millions d'euros déposés par la société (4.47
millions d'euros de marge et 0.5~million d'euros de garantie), 6.21~millions d'euros de crédit accordés par la banque de l'acheteur, et
248\,380~euros par an de carry. \textbf{20.} Collatéral, crédit et conformité avant de prendre tout risque de prix.
\end{solution}

\begin{solution}{iq:m3:getting-access-commodities-and-power:1}
Un compte virtuel auprès de l'opérateur du réseau qui agrège les injections, les soutirages et les transactions d'un portefeuille, et qui est
réglé pour ses déséquilibres. Coûts~: garanties ou sûretés, systèmes de programmation et personnel travaillant sous contrainte de délais,
coûts de déséquilibre, et frais.
\iqlookfor{coûts opérationnels et financiers, et pas seulement la définition.}
\end{solution}

\begin{solution}{iq:m3:getting-access-commodities-and-power:2}
La banque de l'acheteur s'engage à payer contre des documents conformes~: le risque de crédit du vendeur passe de l'acheteur
à la banque, et le paiement ne dépend plus de la coopération de l'acheteur une fois les documents présentés.
\iqlookfor{paiement documentaire et substitution du crédit.}
\end{solution}

\begin{solution}{iq:m3:getting-access-commodities-and-power:3}
Contrats à terme~: marge de variation quotidienne en espèces et marge initiale qui augmente avec la volatilité, de sorte que les fortes variations épuisent
les liquidités avant que la position physique ne rapporte. Forwards de gré à gré~: collatéral sous accords de garantie (ou aucun),
plus de risque de crédit et moins de ponction sur les liquidités. La crise énergétique de 2022 a montré le risque de trésorerie de la voie des contrats à terme.
\iqlookfor{marge en espèces face au risque de crédit.}
\end{solution}

\begin{solution}{iq:m3:getting-access-commodities-and-power:4}
Une relation de compensation qui livrera et recevra des warrants, un entrepôt ou un compte de warrants, des espèces pour payer
le métal qu'elle prend en livraison, et un plan de prêt si elle détient une position dominante.
\iqlookfor{la chaîne de règlement-livraison physique derrière le contrat.}
\end{solution}

\begin{solution}{iq:m3:getting-access-commodities-and-power:5}
Enregistrement \omterm{def:m3:getting-access-commodities-and-power:remit}{REMIT} (et code ACER) auprès d'un régulateur national, dispositions de déclaration des transactions, vérification des
besoins d'autorisation au titre de MiFID II (l'\omterm{def:m3:getting-access-commodities-and-power:aae}{exemption d'activité accessoire} convient rarement à un fonds), accords de bourse et de
compensation, et, pour le négoce physique, contrats de \omterm{def:m3:getting-access-commodities-and-power:group}{groupe d'équilibrage} et de réseau.
\iqlookfor{\omterm{def:m3:getting-access-commodities-and-power:remit}{REMIT}, statut au titre de MiFID II, et la couche physique.}
\end{solution}

\begin{solution}{iq:m3:getting-access-commodities-and-power:6}
Un entrepôt de positions par groupe, zone et MTU, alimenté par toute la capture des transactions~; des données de prévision~; un module de construction des nominations qui
équilibre et apparie par contrepartie~; validation par rapport aux délais et aux formats de l'opérateur~; soumission avec
suivi des accusés de réception et renvoi automatique après les transactions infrajournalières~; rapprochement avec les
programmes appariés de l'opérateur et avec les règlements des déséquilibres.
\iqlookfor{équilibrage, appariement, délais et rapprochement.}
\end{solution}
